% !TEX root = ../licenta.tex
\chapter{Theoretical Foundations}
\label{ch:theoretical_foundations}

This chapter provides the comprehensive mathematical proofs, formal logic, and statistical derivations for the time-series models and structural break detection algorithms utilized natively within the application's Node.js execution engine. A rigorous understanding of these econometric principles is necessary to comprehend how the software manipulates raw data structures into reliable predictive insights without relying on external statistical software.

\section{Time Series Fundamentals and Stationarity}
A time series is defined mathematically as a sequence of random variables, $\{Y_t : t \in \mathbb{T}\}$, observed at discrete, equally spaced time intervals. The fundamental premise of linear time series forecasting is that the structural dependencies observed in the past will persist into the future. 

To formalize this, we rely on Wold's Representation Theorem \cite{Wold1938}, which states that any covariance-stationary time series can be written as the sum of two orthogonal processes: a deterministic process (which can be perfectly predicted from past values) and a stochastic moving average process of infinite order.

For a linear model to be statistically valid under Wold's theorem, the time series must generally exhibit weak (or second-order) stationarity. A stochastic process $Y_t$ is weakly stationary if it satisfies three strict mathematical conditions:
\begin{enumerate}
    \item \textbf{Constant Mean:} The expected value of the process is constant over time. For all $t$, $\mathbb{E}[Y_t] = \mu$. This implies the series does not exhibit a long-term upward or downward trend.
    \item \textbf{Constant Variance:} The variance is finite and constant over time. For all $t$, $\text{Var}(Y_t) = \mathbb{E}[(Y_t - \mu)^2] = \sigma^2 < \infty$. This implies the volatility of the series does not increase or decrease.
    \item \textbf{Autocovariance depends only on lag:} The covariance between any two observations depends strictly on the time lag $k$ between them, not on the actual chronological time $t$. The autocovariance function is defined as:
    \begin{equation}
        \gamma(k) = \text{Cov}(Y_t, Y_{t-k}) = \mathbb{E}[(Y_t - \mu)(Y_{t-k} - \mu)]
    \end{equation}
\end{enumerate}

Consequently, the Autocorrelation Function (ACF), denoted as $\rho(k)$, is the normalized autocovariance used to identify trailing dependencies \cite{Hamilton1994}:
\begin{equation}
    \rho(k) = \frac{\gamma(k)}{\gamma(0)} = \frac{\text{Cov}(Y_t, Y_{t-k})}{\text{Var}(Y_t)}
\end{equation}
Furthermore, the Partial Autocorrelation Function (PACF) isolates the direct correlation between $Y_t$ and $Y_{t-k}$ by mathematically removing the linear dependence on the intermediate lags $Y_{t-1}, Y_{t-2}, \dots, Y_{t-k+1}$. Both the ACF and PACF are critical for diagnosing the mathematical structure of the data before modeling.

\section{Moving Average Models}
The most fundamental approach to smoothing volatile data, extracting underlying signals, and establishing a baseline forecast is the Moving Average. The application implements both Simple and Weighted Moving Averages to accommodate different temporal behaviors.

\subsection{Simple Moving Average (SMA)}
The Simple Moving Average of order $k$, denoted $SMA(k)$, computes the unweighted arithmetic mean of the previous $k$ data points. For a time series $y_t$, the smoothed value at time $t$, which also serves as the na\"ive forecast $\hat{y}_{t+1}$, is calculated as:
\begin{equation}
    \hat{y}_{t+1} = \frac{1}{k} \sum_{i=0}^{k-1} y_{t-i} = \frac{y_t + y_{t-1} + y_{t-2} + \dots + y_{t-k+1}}{k}
\end{equation}

To understand the smoothing effect mathematically, consider the variance of the SMA. If we assume the original series $y_t$ consists of a constant signal $\mu$ plus independent identically distributed (i.i.d.) white noise $\epsilon_t$ with variance $\sigma^2$ (i.e., $y_t = \mu + \epsilon_t$), the variance of the SMA is:
\begin{equation}
    \text{Var}(\hat{y}_{t+1}) = \text{Var}\left( \frac{1}{k} \sum_{i=0}^{k-1} y_{t-i} \right) = \frac{1}{k^2} \sum_{i=0}^{k-1} \text{Var}(y_{t-i}) = \frac{1}{k^2} (k \sigma^2) = \frac{\sigma^2}{k}
\end{equation}
This proof demonstrates that the SMA successfully reduces the variance (noise) of the original series by a factor of $k$. However, the primary limitation of the SMA is its equal weighting and the resulting mathematical lag. Because every point in the window is weighted equally at $\frac{1}{k}$, the "center of mass" of the SMA lags the actual data by precisely $\frac{k-1}{2}$ periods. This makes the SMA notoriously slow to respond to structural breaks or rapid trend reversals.

\subsection{Weighted Moving Average (WMA)}
To mathematically mitigate the lagging effect of the SMA, the Weighted Moving Average (WMA) assigns predefined, decreasing weights to historical data points, ensuring recent data heavily influences the forecast while distant data decays. Let $w_i$ represent the weight applied to the observation $i$ periods ago, where the sum of all weights strictly equals 1: $\sum_{i=1}^{k} w_i = 1$. The formula is:
\begin{equation}
    \hat{y}_{t+1} = \sum_{i=1}^{k} w_i y_{t-i+1}
\end{equation}
In a linearly weighted moving average, the weights decrease arithmetically. The most recent observation $y_t$ receives a weight of $k$, the previous $y_{t-1}$ receives $k-1$, down to a weight of $1$ for $y_{t-k+1}$. To ensure the weights sum to 1, each term is divided by the sum of the integers from 1 to $k$, which is $\frac{k(k+1)}{2}$. Thus, the specific weight $w_i$ is defined as:
\begin{equation}
    w_i = \frac{k - i + 1}{\frac{k(k+1)}{2}} = \frac{2(k - i + 1)}{k(k+1)}
\end{equation}
While the WMA is an improvement over the SMA in tracking active trends with less lag, defining the arbitrary window $k$ remains a subjective constraint.

\section{Exponential Smoothing Methodologies}
Exponential smoothing fundamentally resolves the arbitrary window limitation of moving averages by applying a continuous, geometrically decaying weight to all past observations \cite{Hyndman2008}. This ensures that the most recent data dictates the immediate trajectory of the forecast, while still maintaining the historical context of the entire time series dating back to the genesis point $t=0$.

\subsection{Simple Exponential Smoothing (SES)}
Simple Exponential Smoothing (SES) is utilized when the data, denoted as $y_t$, exhibits a relatively stationary mean with no distinct trend or seasonal pattern. The foundational equation for SES is defined recursively as:
\begin{equation}
    S_t = \alpha y_t + (1 - \alpha) S_{t-1}
\end{equation}
where:
\begin{itemize}
    \item $S_t$ is the smoothed statistic, representing the level of the series at time $t$, which serves as the flat forecast for all future periods: $\hat{y}_{t+h|t} = S_t$.
    \item $y_t$ is the actual observation at time $t$.
    \item $\alpha$ is the smoothing parameter, strictly bounded such that $0 \le \alpha \le 1$.
\end{itemize}

The exponential nature of this model is proven through recursive substitution. If we substitute the previous state $S_{t-1} = \alpha y_{t-1} + (1 - \alpha)S_{t-2}$ into the main equation, we get:
\begin{equation}
    S_t = \alpha y_t + (1 - \alpha)[\alpha y_{t-1} + (1 - \alpha)S_{t-2}] = \alpha y_t + \alpha(1-\alpha)y_{t-1} + (1-\alpha)^2 S_{t-2}
\end{equation}
Continuing this substitution backward to the first observation yields the infinite expansion:
\begin{equation}
    S_t = \alpha \sum_{i=0}^{t-1} (1-\alpha)^i y_{t-i} + (1-\alpha)^t S_0
\end{equation}
This mathematically proves that the weight assigned to past observations decays exponentially by a constant factor of $(1-\alpha)$ as the data ages. An $\alpha$ value approaching 1 heavily discounts older data (reacting rapidly to shocks), while an $\alpha$ value approaching 0 treats all historical data nearly equally (creating a highly smoothed, inert forecast). In practice, optimal initial states ($S_0$) and $\alpha$ are found by minimizing the Sum of Squared Errors (SSE) over the training data.

\subsection{Holt's Linear Trend Model}
When an economic time series possesses an underlying upward or downward trajectory, SES fails because its flat forecasts continuously lag behind the trend. To resolve this, Holt extended the model to explicitly estimate the slope of the trend over time.

The dual recursive equations governing Holt's method are defined mathematically as the Level equation ($L_t$) and the Trend equation ($T_t$):
\begin{align}
    L_t &= \alpha y_t + (1 - \alpha)(L_{t-1} + T_{t-1}) \\
    T_t &= \beta (L_t - L_{t-1}) + (1 - \beta)T_{t-1}
\end{align}
where $\alpha$ is the level smoothing parameter and $\beta$ is the trend smoothing parameter ($0 \le \alpha, \beta \le 1$). 

The level equation updates the current base value by taking a weighted average of the actual observation $y_t$ and the previous level plus the expected trend $(L_{t-1} + T_{t-1})$. The trend equation updates the slope by taking a weighted average of the current perceived change in level $(L_t - L_{t-1})$ and the previously estimated slope $T_{t-1}$.

The $h$-step-ahead forecast is generated by linearly projecting the estimated trend $T_t$ outward from the most recent estimated level $L_t$:
\begin{equation}
    \hat{y}_{t+h|t} = L_t + h T_t
\end{equation}

\subsection{Holt-Winters Additive Method}
To accurately forecast data exhibiting both a linear trend and repeating cyclical variations (seasonality), the Holt-Winters method incorporates a third seasonal parameter. In the context of this architecture, the \textit{additive} formulation is utilized, which assumes that the absolute amplitude of the seasonal fluctuations remains relatively constant regardless of the underlying level of the series.

The system state is updated iteratively across three equations: Level ($L_t$), Trend ($T_t$), and Seasonality ($S_t$):
\begin{align}
    L_t &= \alpha (y_t - S_{t-m}) + (1 - \alpha)(L_{t-1} + T_{t-1}) \\
    T_t &= \beta (L_t - L_{t-1}) + (1 - \beta)T_{t-1} \\
    S_t &= \gamma (y_t - L_{t-1} - T_{t-1}) + (1 - \gamma)S_{t-m}
\end{align}
Here, $\gamma$ is the seasonal smoothing parameter, and $m$ represents the period of the seasonality (e.g., $m=12$ for monthly data with annual cycles, or $m=4$ for quarterly GDP data). 

In the Level equation, the data $y_t$ is dynamically "de-seasonalized" by subtracting the seasonal index from the previous cycle ($S_{t-m}$). In the Seasonal equation, the current seasonal index is estimated by subtracting the current expected level $(L_{t-1} + T_{t-1})$ from the raw data $y_t$, isolating the seasonal variation.

The $h$-step-ahead forecast equation projects the trend, and then adds the appropriate seasonal index isolated from the previous cycle:
\begin{equation}
    \hat{y}_{t+h|t} = L_t + h T_t + S_{t-m+h_m^+}
\end{equation}
where the mathematical notation $h_m^+ = [(h-1) \mod m] + 1$ guarantees that the algorithm accurately cycles back to extract the correct corresponding seasonal multiplier regardless of how far into the future the forecast extends.


\section{Autoregressive Integrated Moving Average (ARIMA)}
The ARIMA$(p, d, q)$ framework, formally synthesized in the Box-Jenkins methodology \cite{BoxJenkins1970}, is one of the most robust and theoretically sound approaches for forecasting stationary and non-stationary linear time series. It structurally combines Autoregressive $AR(p)$ dynamics, differencing $I(d)$ for non-stationary integration, and Moving Average $MA(q)$ error smoothing into a unified predictive equation.

The Box-Jenkins methodology outlines a rigorous three-step process:
\begin{enumerate}
    \item \textbf{Identification:} Determining the parameters $p, d, q$ by analyzing the ACF and PACF plots to identify unit roots and lag dependencies.
    \item \textbf{Estimation:} Using Ordinary Least Squares (OLS) and the Yule-Walker equations to rapidly estimate the exact coefficients for the autoregressive and moving average terms, rather than computationally heavy Maximum Likelihood Estimation (MLE).
    \item \textbf{Diagnostic Checking:} Ensuring the model residuals resemble white noise (via the Ljung-Box test), confirming that all systematic information has been successfully extracted by the model.
\end{enumerate}

\subsection{The Lag Operator and Autoregressive Processes}
To formalize the mathematical derivations of the estimation phase, let $L$ represent the lag (or backshift) operator, defined such that applying $L$ to an observation shifts it back in time by one period: $L y_t = y_{t-1}$. By continuous application, $L^k y_t = y_{t-k}$.

An Autoregressive process of order $p$, denoted $AR(p)$, models the current observation $y_t$ as a linear combination of its $p$ past observations plus a completely stochastic white noise error term $\epsilon_t$:
\begin{equation}
    y_t = c + \phi_1 y_{t-1} + \phi_2 y_{t-2} + \dots + \phi_p y_{t-p} + \epsilon_t
\end{equation}
Utilizing the lag operator, we can move all $y$ terms to the left side and factor out $y_t$:
\begin{equation}
    \left(1 - \phi_1 L - \phi_2 L^2 - \dots - \phi_p L^p\right) y_t = c + \epsilon_t 
\end{equation}
\begin{equation}
    \Phi(L) y_t = c + \epsilon_t
\end{equation}
where $\Phi(L)$ is the autoregressive characteristic polynomial of order $p$. The Yule-Walker equations can be derived from this formulation by multiplying by $y_{t-k}$ and taking the mathematical expectation, allowing for the precise estimation of the $\phi$ parameters based on sample autocorrelations.

\subsection{Moving Average Processes}
Conversely, a Moving Average process of order $q$, denoted $MA(q)$, models the current observation as the mean $\mu$ plus a linear combination of current and $q$ past white noise error terms (historical forecasting shocks):
\begin{equation}
    y_t = \mu + \epsilon_t + \theta_1 \epsilon_{t-1} + \dots + \theta_q \epsilon_{t-q}
\end{equation}
In polynomial lag notation, this is expressed as:
\begin{equation}
    y_t = \mu + \left(1 + \theta_1 L + \theta_2 L^2 + \dots + \theta_q L^q\right) \epsilon_t \implies y_t = \mu + \Theta(L) \epsilon_t
\end{equation}

\subsection{Integration and the Complete ARIMA Model}
If the original economic time series is non-stationary (e.g., exhibiting stochastic drift or random walks), it must be differenced $d$ times to achieve stationarity. The first difference is calculated as $\Delta y_t = y_t - y_{t-1} = (1-L)y_t$. The $d$-th difference is denoted as $(1-L)^d y_t$. 

Combining the Autoregressive polynomial, the differencing integration operator, and the Moving Average polynomial yields the formal, canonical definition of the $ARIMA(p, d, q)$ stochastic process:
\begin{equation}
    \Phi(L) (1-L)^d y_t = c + \Theta(L) \epsilon_t
\end{equation}
where $\epsilon_t \sim WN(0, \sigma^2)$ is strictly assumed to be a white noise process characterized by a mean of zero, constant variance $\sigma^2$, and zero covariance between distinct periods ($\text{Cov}(\epsilon_t, \epsilon_{t-k}) = 0$ for all $k \neq 0$).

\subsection{Stationarity and Invertibility Roots}
The theoretical viability of an ARIMA model depends explicitly on the complex roots of its characteristic polynomials. For the $AR(p)$ component of the model to represent a strictly stationary, non-explosive process, the roots of the characteristic equation $\Phi(z) = 0$ must lie strictly outside the unit circle in the complex plane. Mathematically:
\begin{equation}
    1 - \phi_1 z - \phi_2 z^2 - \dots - \phi_p z^p = 0
\end{equation}
For any complex root $z$ satisfying this polynomial, it is an absolute mathematical requirement that the modulus $|z| > 1$. If a root lies exactly on the unit circle ($|z| = 1$), the process exhibits a "unit root," indicating severe non-stationarity, and requires a higher degree of differencing (an increase in the $d$ hyper-parameter).

Correspondingly, for the $MA(q)$ process to be "invertible"—meaning the unobservable error term $\epsilon_t$ can be mathematically inverted and rewritten as an infinite converging sum of past observable data $y_t$—all roots of the moving average polynomial $\Theta(z) = 1 + \theta_1 z + \dots + \theta_q z^q = 0$ must also fall strictly outside the unit circle ($|z| > 1$). Without invertibility, the moving average parameters cannot be uniquely identified from the autocorrelation function of the sample data.


\section{Change Point Detection (CPD) Mathematics}
As repeatedly established, economic time series are highly susceptible to structural breaks caused by systemic shocks. Fitting a single global model across these breaks violates stationarity assumptions and ruins forecast accuracy. The mathematical process of identifying these shifts is Change Point Detection (CPD).

Formally, a time series $y_{1:n} = (y_1, y_2, \dots, y_n)$ contains $m$ change points if it can be split into $m+1$ segments, where the statistical properties within each segment are constant, but the properties between adjacent segments differ. Let the positions of these change points be denoted by $\tau = (\tau_1, \tau_2, \dots, \tau_m)$, where $0 = \tau_0 < \tau_1 < \tau_2 < \dots < \tau_m < \tau_{m+1} = n$.

The goal of a multiple change point algorithm is to jointly estimate both the number of change points $m$ and their exact temporal locations $\tau_{1:m}$. This is formulated as an optimization problem minimizing a penalized objective function:
\begin{equation}
    \min_{m, \tau_{1:m}} \left\{ \sum_{i=1}^{m+1} \left[ \mathcal{C}(y_{(\tau_{i-1}+1):\tau_i}) \right] + \beta f(m) \right\}
\end{equation}
where:
\begin{itemize}
    \item $\mathcal{C}(\cdot)$ is a predefined cost function measuring the fit of the data within a specific segment. Common cost functions include the negative log-likelihood or the residual sum of squares (variance).
    \item $\beta f(m)$ is a penalty term designed to guard against overfitting. Without a penalty ($\beta = 0$), the algorithm would trivially place a change point at every single data coordinate to drive the cost to zero. 
\end{itemize}

\section{The Pruned Exact Linear Time (PELT) Algorithm}
To solve the CPD optimization problem, one could theoretically evaluate every possible combination of breakpoints. Early algorithms like Binary Segmentation attempted to solve this heuristically by finding a single change point, splitting the data in two, and recursively searching each half. However, this is an approximate method and does not guarantee a globally optimal solution. Conversely, evaluating all combinations via an exact brute-force search is computationally unfeasible, scaling at an exponential rate of $\mathcal{O}(2^n)$. 

The Optimal Partitioning (OP) method, based on dynamic programming, reduces this complexity to $\mathcal{O}(n^2)$. It defines $F(t)$ as the absolute minimum optimal cost for data up to time $t$:
\begin{equation}
    F(t) = \min_{0 \le \tau < t} \left[ F(\tau) + \mathcal{C}(y_{(\tau+1):t}) + \beta \right]
\end{equation}
While $\mathcal{O}(n^2)$ is significantly better, it is still too slow for instantaneous, interactive web environments processing tens of thousands of data points. The solution implemented in this application is the Pruned Exact Linear Time (PELT) algorithm, formalized by Killick, Fearnhead, and Eckley \cite{Killick2012}.

\subsection{Mathematical Proof of Pruning Logic}
The brilliance of the PELT algorithm lies in its mathematical pruning condition, which safely discards candidate change points that can mathematically never be the optimal choice. This reduces the complexity to exact linear time, $\mathcal{O}(n)$.

Theorem 3.1 of Killick et al. (2012) dictates the pruning rule. Assume there exists a constant $K$ such that the cost function satisfies the following inequality for all $\tau < t < T$:
\begin{equation}
    \mathcal{C}(y_{(\tau+1):t}) + \mathcal{C}(y_{(t+1):T}) + K \le \mathcal{C}(y_{(\tau+1):T})
\end{equation}
This condition simply requires that introducing a new change point at time $t$ reduces the total cost of the segment by at least $K$. For the negative log-likelihood cost function, $K=0$. 

The theorem mathematically proves that if, at any given time $t$, there exists a previous time point $\tau$ that satisfies the inequality:
\begin{equation}
    F(\tau) + \mathcal{C}(y_{(\tau+1):t}) + K \ge F(t)
\end{equation}
then the point $\tau$ can never be the optimal most recent change point for any future time $T > t$. 

Because this condition guarantees that $\tau$ is permanently suboptimal, the temporal index $\tau$ is forcefully pruned from the set of candidate change points. As the algorithm progresses forward through the time series, the list of potential previous change points remains incredibly small, resulting in $\mathcal{O}(n)$ linear computational performance. This extreme algorithmic efficiency is what allows the Node.js V8 application architecture to calculate structural breaks natively and synchronously without UI thread starvation, completely bypassing the need for computationally heavy external Python sub-processes.

\section{Derivation of the Bayesian Information Criterion (BIC)}
To ensure the PELT algorithm accurately penalizes the addition of new change points, an optimal penalty constant $\beta$ must be calculated. The application defaults to the Bayesian Information Criterion (BIC), which is mathematically derived to approximate the true posterior probability of a model, aggressively penalizing complexity to find the true structural breaks \cite{Schwarz1978}.

The derivation begins with the marginal likelihood of the data $y$ given a specific model structure $m$ (representing a specific set of change points):
\begin{equation}
    p(y | m) = \int p(y | \theta, m) p(\theta | m) d\theta
\end{equation}
where $\theta$ represents the parameters within the segments, $p(y | \theta, m)$ is the likelihood function, and $p(\theta | m)$ is the prior distribution of the parameters. 

This integral is highly complex. However, applying Laplace's method for integral approximation, we assume that for large sample sizes, the posterior distribution $p(\theta | y, m)$ is highly peaked around its mode, which corresponds to the maximum likelihood estimator $\hat{\theta}$. 

Taking a second-order Taylor series expansion of the log-posterior around the maximum $\hat{\theta}$, we obtain:
\begin{equation}
    \log p(y | m) \approx \log p(y | \hat{\theta}, m) + \log p(\hat{\theta} | m) + \frac{k}{2} \log(2\pi) - \frac{1}{2} \log |H|
\end{equation}
where $k$ is the number of free parameters in model $m$, and $H$ is the Hessian matrix (the matrix of second derivatives) of the negative log-likelihood function evaluated at $\hat{\theta}$. 

As the sample size $n \to \infty$, the terms involving the prior distribution $p(\hat{\theta} | m)$ and the constant $\frac{k}{2} \log(2\pi)$ do not grow with $n$ and thus become asymptotically negligible. The determinant of the Hessian matrix, however, grows at a rate directly proportional to $n^k$. Therefore, we can approximate $\log |H| \approx k \log n$.

Substituting this back into the expansion, we achieve the large-sample asymptotic approximation for the log marginal likelihood:
\begin{equation}
    \log p(y | m) \approx \log p(y | \hat{\theta}, m) - \frac{k}{2} \log n
\end{equation}
To convert this into a deviance metric (where smaller values indicate a better model fit), the entire equation is multiplied by $-2$:
\begin{equation}
    -2 \log p(y | m) \approx -2 \log p(y | \hat{\theta}, m) + k \log n
\end{equation}
This final expression corresponds precisely to the formal BIC equation:
\begin{equation}
    \text{BIC} = -2 \log(\hat{L}) + k \log n
\end{equation}
where $\hat{L} = p(y | \hat{\theta}, m)$ represents the maximized likelihood of the model. In the context of the PELT algorithm, utilizing the BIC implies setting the penalty parameter $\beta = \log n$. This mathematical derivation proves that the BIC scales the penalty directly logarithmically with the number of data points, fiercely protecting the forecast from identifying false-positive structural shifts caused by random noise.

\section{Segmented Forecasting Mechanics and RMSE Validation}
Once the PELT algorithm successfully identifies a valid change point $\tau_{break}$, the architecture relies on the mechanical execution of a Segmented Forecast.

The logic dictates that the time series $y_{1:n}$ must be completely severed at $\tau_{break}$. All data points in the pre-shock regime, defined as $Y_{pre} = \{y_t : 1 \le t < \tau_{break}\}$, are strictly eliminated from the predictive arrays. The newly truncated dataset, $Y_{post} = \{y_t : t \ge \tau_{break}\}$, is designated as the new training matrix. This matrix is passed independently into the ARIMA or Holt-Winters execution blocks. 

Because the non-stationary volatility of $Y_{pre}$ is permanently removed, the optimization functions calculating the internal parameters converge on entirely different, highly localized constants. This entirely eliminates the heavy lag bias that traditional global models suffer from during macroeconomic recoveries.

To mathematically validate this architectural strategy to the end-user, the backend calculates the Root Mean Squared Error (RMSE) for both the global model and the segmented model. The RMSE is defined as the square root of the average of squared differences between prediction and actual observation:
\begin{equation}
    \text{RMSE} = \sqrt{ \frac{1}{N_{post}} \sum_{t=\tau_{break}}^{n} (y_t - \hat{y}_t)^2 }
\end{equation}
where $N_{post}$ is the number of observations in the post-break regime. Importantly, to ensure a statistically fair comparison, the RMSE for the global model is evaluated \textit{only} over the same out-of-sample post-break timeline $t \ge \tau_{break}$, measuring precisely how poorly the pre-shock data polluted the current recovery trend. 

While Mean Absolute Error (MAE) could also be used, RMSE is chosen explicitly because squaring the residuals penalizes large errors exponentially, which is highly preferred when modeling chaotic financial markets. The ultimate measure of architectural success is the accuracy gain delta, calculated as:
\begin{equation}
    \Delta_{accuracy} = \left( \frac{\text{RMSE}_{global} - \text{RMSE}_{seg}}{\text{RMSE}_{global}} \right) \times 100\%
\end{equation}

\section{Deterministic AI Prompt Engineering and Execution}
The final theoretical component of the application bridges the gap between raw statistical output and natural language business analytics. To achieve this, the architecture integrates the Google Gemini Large Language Model (LLM). 

As previously established, standard LLMs inherently utilize probabilistic, next-token prediction algorithms governed by self-attention mechanisms \cite{Vaswani2017}. While excellent for generating natural language, this probabilistic nature makes them highly prone to mathematical hallucination when asked to independently analyze numerical data. To completely restrict the AI and force a deterministic outcome, the backend server utilizes a rigid "Heuristic Prompt Injection" methodology. 

Rather than sending the raw data array to the LLM via an API call, the Node.js backend executes the mathematical formulas described above, calculates the exact $\Delta_{accuracy}$, and explicitly injects these constants into a rigid JSON template. The system prompts the LLM with instructions such as: 

\begin{quote}
\textit{"A structural break occurred at date X. The global forecast produced an error rate of $Y$. By segmenting the data and deleting obsolete history, the new error rate is $Z$, yielding a proven mathematical accuracy improvement of $W\%$. Acting as a strict econometrician, explain the semantic, macroeconomic events that occurred on date X that justify why discarding the historical data improved the mathematical fit."}
\end{quote}

Furthermore, the API call to the LLM is executed with the `temperature` parameter explicitly set to $0.0$. In LLM architecture, temperature governs the randomness of the softmax distribution across the output vocabulary. By forcing the temperature to zero, the LLM is restricted to purely greedy decoding, selecting only the most probable token at every step. 

By denying the LLM the ability to calculate the numbers itself, forcing a temperature of zero, and compelling it to narrate pre-calculated, mathematically proven econometric variables, the architecture guarantees a deterministic, hallucination-free analytical output. This effectively hybridizes the rigorous statistical foundations of time-series mathematics with the vast, contextual intelligence of modern generative AI.
